\chapter{Introdução}


A discussão sobre a eficiência energética de tecnologias da informação e
comunicação (TIC) sempre teve como foco a melhoria dos componentes de
\textit{hardware}. \citeonline{koomeys-law} observaram uma tendência de
melhoria exponencial desses dispositivos, com duplicação aproximada da
eficiência energética a cada um ano e meio, tendência conhecida popularmente
como Lei de Koomey. Essa tendência de crescimento exponencial se confirma mais
recentemente, ainda que em um ritmo mais lento do que o previsto anteriormente
\cite{koomeys-law-revisited}. O sucesso em otimizar a eficiência do
\textit{hardware} oculta, porém, outro aspecto das TIC que recebe menos
atenção: o \textit{software} \cite{is-software-green}. Os esforços para
melhorar a eficiência das TIC devem ser compartilhados por desenvolvedores de
\textit{hardware} e \textit{software} se quisermos potencializar essa economia
de energia, o que se torna cada vez mais relevante com o crescimento dos
dispositivos móveis.

Apesar da predominância do \textit{hardware}, o interesse dos desenvolvedores
em assuntos relacionados ao consumo de energia em \textit{software} é crescente
\cite{mining-questions}. Isso evidencia a relevância do tema. Além desse
crescimento, \citeonline{mining-questions} também mostram um engajamento dos
desenvolvedores maior que a média quando comparado com outros temas. Por outro
lado, o mesmo estudo mostra que o entendimento sobre o assunto ainda era
limitado, apoiando-se em soluções genéricas e em discussões pouco aprofundadas.
Além disso, de forma mais abrangente, o campo de \textit{green computing}
(computação verde) atrai cada vez mais atenção dos pesquisadores, considerando
que o número de publicações ligadas ao tema aumentou de forma consistente entre
2007 e 2022 \cite{green-computing-past-present-future,
green-computing-literature-survey}. O aumento da relevância do tema ressalta a
importância de fornecer a desenvolvedores e pesquisadores métricas precisas
relacionadas à energia, para que possam analisar os impactos de suas decisões
na eficiência energética das aplicações.

O uso de \textit{benchmarks} para avaliação de \textit{software} é parte
importante do seu processo de desenvolvimento. Porém, métricas relacionadas ao
consumo de energia e à eficiência energética dessas aplicações raramente são
consideradas na elaboração desses \textit{benchmarks}. No ecossistema Java, as
ferramentas mais utilizadas para esse tipo de análise, como os conjuntos de
\textit{benchmarks} Renaissance \cite{renaissance}, DaCapo
\cite{dacapochopin} e SPEC \cite{spec}, não incluem métricas relacionadas ao
desempenho sob a perspectiva do consumo de energia. Isso representa uma
limitação importante, que pode levar desenvolvedores e pesquisadores a tomar
decisões menos bem informadas, resultando em custos energéticos mais altos que
passam despercebidos. A primeira contribuição deste trabalho é o
desenvolvimento de uma extensão ao pacote de \textit{benchmarks} DaCapo
Chopin, adicionando uma forma de medir o consumo de energia.

Devido à falta de medidas relacionadas à energia nessas ferramentas, a
avaliação da eficiência energética das aplicações é muitas vezes inferida a
partir do tempo de execução. Ainda que haja correlação entre essas duas
métricas, consumo de energia e tempo de execução não estão perfeitamente
relacionados \cite{programming_lang_energy}. Isso justifica a importância de
testá-las e analisá-las separadamente. A partir da extensão desenvolvida e dos
dados disponibilizados por \citeonline{dacapochopin}, foi realizada uma análise
do consumo de energia e da eficiência energética para 16 cargas de trabalho da
ferramenta, sob diferentes configurações da Java Virtual Machine (JVM) e do
\textit{hardware} utilizado. Com base nessas métricas, este trabalho tem como
objetivo compreender o comportamento do consumo de energia da JVM para os
programas avaliados.